\section{How to use \texttt{MONKES}}


\subsection{Magnetic configuration input}\label{sec:Magnetic_configuration_input}


{\MONKES} can read the magnetic configuration of a particular flux surface from different types of files.
\subsubsection{From {\BOOZERXFORM} output file}
{\MONKES} can read the magnetic confituration of an output (binary) file from {\BOOZERXFORM}. The file \textbf{must} be named \texttt{boozmn.nc}. It is advisable to make a soft link instead of renaming the original {\BOOZERXFORM} output file. 

In an output file from {\BOOZERXFORM} there is information about the equilibrium of many flux surfaces. If not specified, {\MONKES} will perform its calculations at the flux surface labelled by $s={\psi/\psi_{\text{lcfs}}}=0.25$ where $2\pi\psi$ is the toroidal flux and $\psi_{\text{lcfs}}$ the label of the last closed flux surface. A different flux surface can be specified in the input file \texttt{monkes\_input.surface}. An example on how to fill the input file \texttt{monkes\_input.surface} is given in listing \ref{lst:monkes_input.surface}. 


\listings{../../examples/W7X-EIM/boozer_xform_input/monkes_input.surface}{&surface}{/}{Example of \texttt{monkes\_input.surface} for specifying the surface $s={\psi/\psi_{\text{lcfs}}}=0.2$}{lst:monkes_input.surface}

\subsubsection{From {\DKES} input file}
The magnetic configuration can also be specified using a {\DKES} input file. This input file contains the information about a single flux surface and therefore $s$ does not need to be (explicitly) specified. For running {\MONKES} using a {\DKES} input file, it \textbf{must} be renamed as \texttt{ddkes2.data}. 

%\subsubsection{From {\DESC} input file}

\subsection{Monoenergetic database input}\FloatBarrier
In addition to the magnetic configuration (see section \ref{sec:Magnetic_configuration_input}), in order to run {\MONKES} for calculating the monoenergetic coefficients $\Dij{ij}$, it is necessary to specify 5 additional parameters. The required parameters are listed in table \ref{tab:monkes_input.parameters}.

\begin{table}[]
	\centering
	\begin{tabular}{@{}llll@{}}
		\toprule
		Parameter             & Symbol in {\MONKES} paper \cite{Escoto_2024} & {Name in {\MONKES} input file} & Units           \\ \midrule
		Poloidal resolution   & $N_\theta$                                 & \texttt{N\_theta}                                & N/A             \\
		Toroidal resolution   & $N_\zeta$                                 & \texttt{N\_zeta}                                & N/A             \\
		No. Legendre modes    & $N_\xi$                                    & \texttt{N\_xi}                                   & N/A             \\
		Collisionallity       & $\hat{\nu}$                                & \texttt{nu}                                      & $1/\text{m}$    \\
		Radial electric field & $\Er$                                      & \texttt{E\_r}                                    & $\text{V}\cdot \text{m}/\text{s}$ \\ \bottomrule
	\end{tabular}
	\caption{Parameters to be specified in the file \texttt{monkes\_input.parameters} for calculating the monoenergetic coefficients.}
	\label{tab:monkes_input.parameters}
\end{table}

These parameters have to be specified in the input file \texttt{monkes\_input.parameters}. For each combination of $(N_\theta,N_\zeta,N_\xi,\hat{\nu},N_\xi)$ {\MONKES} will compute the monoenergetic coefficients $\Dij{ij}$. An example on how to fill this input file is given in listing \ref{lst:monkes_input.parameters}. Note that in the example shown in listing \ref{lst:monkes_input.parameters} there is a single value for \texttt{N\_theta}, \texttt{N\_zeta} and \texttt{N\_xi} while there are 10 values \texttt{nu} and 2 values of \texttt{E\_r}. This means that {\MONKES} will perform a scan in $\hat{\nu}$ and $\Er$ for fixed resolution $(N_\theta,N_\zeta,N_\xi)=(30,60,160)$. Similarly, one could be interested in doing scans in the different resolutions to check how the monoenergetic coefficients $\Dij{ij}$ converge. This scan is also possible by introducing several values of, say $N_\zeta$, separated by commas. 

\listings{../../examples/W7X-EIM/dkes_input/monkes_input.parameters}{&parameters}{/}{Example of \texttt{monkes\_input.parameters} for calculating a monoenergetic database}{lst:monkes_input.parameters}




\subsection{Monoenergetic database output}

The monoenergetic coefficients $\Dij{ij}$ are written in an output file called \texttt{monkes\_Monoenergetic\_Database.dat}. Instead of writing directly $\Dij{ij}$ which are in meters, in the output file \texttt{monkes\_Monoenergetic\_Database.dat}, the quantities written are $K_{ij} \Dij{ij}$. The monoenergetic coefficients are defined in section 2 of reference \cite{Escoto_2024}. The quantities $K_{ij}$, as defined in Appendix E of \cite{Escoto_2024}, are: $K_{11}=(\dv*{\psi}{r})^{-2}$, $K_{13}=K_{31}=(\dv*{\psi}{r})^{-1}$ and $K_{33}=1$.


\begin{landscape}
	\listings{../../examples/W7X-EIM/dkes_input/monkes_Monoenergetic_Database.dat}{nu}{0.3000000000000000E+00}{Example of \texttt{monkes\_Monoenergetic\_Database.dat} for calculating a monoenergetic database}{lst:monkes_Monoenergetic_Database.dat}
\end{landscape}


\subsection{Example of running {\MONKES} for Wendelstein 7-X EIM geometry}

In the folder location \texttt{\detokenize{examples/W7X-EIM/}} there is an example for running {\MONKES} in Wendelstein 7-X (W7-X) geometry. For both cases, the \texttt{\detokenize{monkes_input.parameters}} file used is the one shown in listing \ref{lst:monkes_input.parameters}. That is, a scan of 10 points in $\hat{\nu}$ for 2 different values of the radial electric field $\Er$.


\subsubsection{Running from {\DKES} input}\label{sec:Running_example_dkes_input}
An example in which the magnetic configuration is specified via a {\DKES} input file can be found at \texttt{\detokenize{examples/W7X-EIM/dkes_input/}}. Once that {\MONKES} has been correctly compiled, this example can be run by simply typing
\begin{verbatim}
	./../../bin/main_monkes.x
\end{verbatim} 
in the command prompt or, if running in a remote platform, by executing the batch file from within this location.

The monoenergetic coefficients, are written in the \texttt{\detokenize{monkes_Monoenergetic_Database.dat}} file. Specifically, the results of this run correspond to the output file example of listing \ref{lst:monkes_Monoenergetic_Database.dat}. In figure \ref{fig:Monoenergetic_database_dkes_input}, the main results of the simulation are plotted, for both values of the electric field $\Er$, against the collisionality $\hat{\nu}$. 
\begin{figure}[H]
	\tikzsetnextfilename{D11_W7X_EIM_dkes_input}
	\begin{tikzpicture}
		\begin{axis}[  
			xmode=log, ymode=log,
			no markers, name=D11W7XEIM, width=0.45\textwidth,xlabel=$\hat{\nu}$ {$[1/\text{m}]$}, ylabel={{$ K_{11}\Dij{11}$ {[$\text{T}^{-2}\cdot\text{m}^{-1}$]} }}, ]
			\addplot[blue] table [skip first n=1, restrict expr to domain={\thisrowno{1}}{0:0},
			x expr=\thisrowno{0},
			y expr=\thisrowno{5}] {../../examples/W7X-EIM/dkes_input/monkes_Monoenergetic_Database.dat}; 
			\addplot[red] table [skip first n=1, restrict expr to domain={\thisrowno{1}}{3e-4:3e-4},
			x expr=\thisrowno{0},
			y expr=\thisrowno{5}] {../../examples/W7X-EIM/dkes_input/monkes_Monoenergetic_Database.dat}; 
			
			\legend{$\Er=0$, $\Er=\num{3e-4}$}
		\end{axis}  
	\end{tikzpicture} 
	\tikzsetnextfilename{D31_W7X_EIM_dkes_input}
	\begin{tikzpicture}
	\begin{axis}[  
		xmode=log,  
		no markers, name=D11W7XEIM, width=0.45\textwidth, xlabel=$\hat{\nu}$ {$[1/\text{m}]$}, ylabel={{$ K_{31}\Dij{31} $ {[$\text{T}^{-1}$]} }}, ]
		\addplot[blue] table [skip first n=1, restrict expr to domain={\thisrowno{1}}{0:0},
		x expr=\thisrowno{0},
		y expr=\thisrowno{6}] {../../examples/W7X-EIM/dkes_input/monkes_Monoenergetic_Database.dat}; 
		\addplot[red] table [skip first n=1, restrict expr to domain={\thisrowno{1}}{3e-4:3e-4},
		x expr=\thisrowno{0},
		y expr=\thisrowno{6}] {../../examples/W7X-EIM/dkes_input/monkes_Monoenergetic_Database.dat}; 
		
		\legend{$\Er=0$, $\Er=\num{3e-4}$}
	\end{axis}  
	\end{tikzpicture} 
	\caption{Radial transport $K_{11}\Dij{11}$ (left) and bootstrap current coefficient $K_{31}\Dij{31}$ (right) obtained from running the example using a {\DKES} input file.}
	\label{fig:Monoenergetic_database_dkes_input}
\end{figure}
	
	
\subsubsection{Running from {\BOOZERXFORM} output}\label{sec:Running_example_boozer_xform_input}
An example in which the magnetic configuration is specified via a {\BOOZERXFORM} output file can be found at \texttt{\detokenize{examples/W7X-EIM/boozer_xform_input/}}. In this folder, the file \texttt{\detokenize{boozmn_wout_w7x_std_v_0.0.txt.nc}} contains the magnetic configuration. As mentioned in section \ref{sec:Magnetic_configuration_input} {\MONKES} expects any {\BOOZERXFORM} output file to be named \texttt{boozmn.nc}. Thus, before running {\MONKES}, the file \texttt{\detokenize{boozmn_wout_w7x_std_v_0.0.txt.nc}} should be renamed to \texttt{boozmn.nc}. The surface $\psi/\psi_{\text{lcfs}}=0.2$ is specified in the \texttt{\detokenize{monkes_input.surface}} file as shown in listing \ref{lst:monkes_input.surface}. This surface also corresponds to the one specified in the {\DKES} input file from section \ref{sec:Running_example_boozer_xform_input}.



Now, provided that {\MONKES} has been correctly compiled, this example can be run by simply typing
\begin{verbatim}
	./../../bin/main_monkes.x
\end{verbatim}
in the command prompt or, if running in a remote platform, by executing the batch file from within this location.

In figure \ref{fig:Monoenergetic_database_boozerxform_output}, the main results of the simulation (written in the \texttt{\detokenize{monkes_Monoenergetic_Database.dat}} file) are plotted. 
\begin{figure}[H]
	\tikzsetnextfilename{D11_W7X_EIM_boozxform_input}
	\begin{tikzpicture}
		\begin{axis}[  
			xmode=log, ymode=log,
			no markers, name=D11W7XEIM,width=0.45\textwidth,xlabel=$\hat{\nu}$ {$[1/\text{m}]$}, ylabel={{$ K_{11}\Dij{11}$ {[$\text{T}^{-2}\cdot\text{m}^{-1}$]} }}, ]
			\addplot[blue] table [skip first n=1, restrict expr to domain={\thisrowno{1}}{0:0},
			x expr=\thisrowno{0},
			y expr=\thisrowno{5}] {../../examples/W7X-EIM/boozer_xform_input/monkes_Monoenergetic_Database.dat}; 
			\addplot[red] table [skip first n=1, restrict expr to domain={\thisrowno{1}}{3e-4:3e-4},
			x expr=\thisrowno{0},
			y expr=\thisrowno{5}] {../../examples/W7X-EIM/boozer_xform_input/monkes_Monoenergetic_Database.dat}; 
			
			\legend{$\Er=0$, $\Er=\num{3e-4}$}
		\end{axis}  
	\end{tikzpicture} 
	\tikzsetnextfilename{D31_W7X_EIM_boozxform_input}
	\begin{tikzpicture}
		\begin{axis}[  
			xmode=log,  
			no markers, name=D11W7XEIM,width=0.45\textwidth,xlabel=$\hat{\nu}$ {$[1/\text{m}]$}, ylabel={{$ K_{31}\Dij{31}$ {[$\text{T}^{-1}$]} }}, ]
			\addplot[blue] table [skip first n=1, restrict expr to domain={\thisrowno{1}}{0:0},
			x expr=\thisrowno{0},
			y expr=\thisrowno{6}] {../../examples/W7X-EIM/boozer_xform_input/monkes_Monoenergetic_Database.dat}; 
			\addplot[red] table [skip first n=1, restrict expr to domain={\thisrowno{1}}{3e-4:3e-4},
			x expr=\thisrowno{0},
			y expr=\thisrowno{6}] {../../examples/W7X-EIM/boozer_xform_input/monkes_Monoenergetic_Database.dat}; 
			
			\legend{$\Er=0$, $\Er=\num{3e-4}$}
		\end{axis}  
	\end{tikzpicture} 
	\caption{Radial transport $K_{11}\Dij{11}$ (left) and bootstrap current coefficient $K_{31}\Dij{31}$ (right) obtained from running the example using a {\BOOZERXFORM} output file.}
	\label{fig:Monoenergetic_database_boozerxform_output}
\end{figure}